\documentclass[main.tex]{subfiles}


\begin{document}

%from pagenumber to up
% \phantomsection 
% \hypertarget{toc}{}

 

 

% номер секции вручную
\setcounter{section}{12
}
\addtocounter{section}{-1} 

\section{Дисперсия света}


Прежде всего нужно сказать, что белый (солнечный) свет является смесью световых пучков разных цветов, а именно~--- смесью волн всего видимого диапазона (от красного до фиолетового).  

На рис.\,\ref{fig:p194} показан \emph{спектр} белого света~--- полоса со всеми цветами радуги, плавно переходящими друг в друга.  

    \begin{figure}[H]
    \centering

    \begin{tikzpicture}[font=\small,            line cap=round,                             >={Stealth[inset=0pt,round,length=3mm, angle'=20]},vec/.style={>={Stealth[inset=0pt,round,length=3.5mm, angle'=20]}}]


\draw[->] (-1,0) -- (13,0) node[below] {$\nu$, Гц};
\node[inner sep=0mm,label={below:$4\cdot10^{14}$}] at (1,0) {\tikz\draw(0,0) -- (0,2mm);};
\node[inner sep=0mm,label={below:$8\cdot10^{14}$}] at (11,0) {\tikz\draw(0,0) -- (0,2mm);};


\node[anchor=south west,inner sep=0] at (1,0.2) {\pgfspectra[width=10cm,gamma=0,begin=740,end=380]};



    \end{tikzpicture}
\caption{Спектр белого света}
\label{fig:p194}
\end{figure}


Каждому цвету соответствует волна определенной частоты в диапазоне видимого света (волны с частотами примерно от $4\cdot10^{14}$~Гц до $8\cdot10^{14}$~Гц). Красный свет имеет наименьшую частоту, а фиолетовый~--- наибольшую. 

В спектре белого света удобно выделять семь цветов~--- красный, оранжевый, желтый, зеленый, голубой, синий, фиолетовый\footnote{Для запоминания цветов можно использовать такое правило: <<каждый охотник желает знать, где сидит фазан>> (первые буквы слов фразы соответствует первым буквам цветов).}. Так цвета перечислены \emph{в порядке возрастания частоты} света (см.~рис.\,\ref{fig:p194}).

\textbf{Дисперсия света}~--- это зависимость преломляемости луча света от его частоты (на данной границе сред).

Пусть узкий луч солнечного света падает на стеклянную призму (рис.\,\ref{fig:p195}). 

    \begin{figure}[H]
    \centering

    \begin{tikzpicture}[font=\small,            line cap=round,                             >={Stealth[inset=0pt,round,length=3mm, angle'=20]},vec/.style={>={Stealth[inset=0pt,round,length=3.5mm, angle'=20]}}]


\fill[lightgray,opacity=.6] (0,0) rectangle (\linewidth,5); 
\coordinate (O) at (.5\linewidth,2.5);

\def\R{4}
\fill[SkyBlue,nearly transparent] (O) ++ (0,{-sqrt(3)/4*\R}) --++ (180:{\R/2}) --++ (60:\R) --++ (-60:\R) -- cycle;
\draw[] (O) ++ (0,{-sqrt(3)/4*\R}) --++ (180:{\R/2})  --++ (60:{\R/2}) coordinate (lt) --++ (60:{\R/2}) coordinate (at) --++ (-60:\R) coordinate (rt) -- cycle;
\path[name path=a] (at) -- (rt);


% color
\def\wc{30}
\draw[white,very thick,vec,->] (lt) ++ ({\wc+180}:2) --++ ({\wc}:1);
\draw[white,very thick] (lt) ++ ({\wc+180}:2) ++ ({\wc}:.9) --++ ({\wc}:1.1);
% \draw[dashed] (lt) ++ ({60+90}:2) --++ ({60+90+180}:4);

\foreach \i/\col in {1,2,3,4,5,6,7}{
\def\gam{1.8}
\path[semithick,name path global=b\i] (lt) --++ ({\wc-5-.7*\i^\gam}:2.5);
% \coordinate[name intersections={of=a and b\i, by=x}] (x\i) at (x);
}

\definecolor{viol}{RGB}{139,0,255}
\definecolor{oran}{RGB}{255,165,0}
\foreach \i/\col in {1/red,2/red!40!yellow,3/yellow,4/Green,5/Cerulean,6/blue,7/viol}{
\def\gam{1.8}
\draw[semithick,\col,name intersections={of=a and b\i, by=x}] (lt) -- (x);
% \coordinate[name intersections={of=a and b\i, by=x}] (x\i) at (x);
}



% \draw[name intersections={of=a and b7, by=x}] (x) --++ (2,0);


% \node[draw] (r) at (O) {\pgfmathparse{sin(90-\wc)/sin(30+\wc-5-.7*1*1))}\pgfmathresult};
% \node[draw] (r) at (O) {\pgfmathparse{sin(90-\wc)/sin(20))}\pgfmathresult};



%after
% \draw[dashed,name intersections={of=a and b1, by=x}] (x) ++ ({\wc}:2) --++ ({\wc+180}:4);
\foreach \i/\col in {1/red,2/red!40!yellow,3/yellow,4/Green,5/Cerulean,6/blue,7/viol}{
\def\gam{1.8}
\def\n{sin(90-\wc)/sin(\wc+\wc-5-.7*\i^\gam)}
\def\al{\wc-(\wc-5-.7*\i^\gam)}
\def\be{asin(sin(\al)*\n)}
\draw[thick,\col,name intersections={of=a and b\i, by=x}] (x) --++ ({\wc-\be}:2);

}



% \node[draw] (r) at (O) {\pgfmathparse{sin(90-\wc)/sin(\wc+\wc-5-.7*1*1)}\pgfmathresult};
% \node[draw] (r) at (O) {\pgfmathparse{\wc-(\wc-5-.7*1*1)}\pgfmathresult};
% \node[draw] (r) at (O) {\pgfmathparse{sin(5)*1.1}\pgfmathresult};








% \def\nr{1.515}
% \def\no{1.516}
% \def\ny{1.517}
% \def\ngr{1.520}
% \def\ng{1.523}
% \def\nb{1.527}
% \def\nv{1.532}

% \node[draw] (r) at (O) {\pgfmathparse{asin(sin(90-\wc)/\nb)}\pgfmathresult};
% \node[below= of r] (o) at (O) {\pgfmathparse{asin(sin(90-\wc)/\nv)}\pgfmathresult};







% \draw (O) circle (2.5);

    
    
    
    
    
    
    \end{tikzpicture}
\caption{Дисперсия света}
\label{fig:p195}
\end{figure}

Можно считать, что на призму падают разноцветные лучи, смешанные в один луч (эту <<смесь>> глаз воспринимает как белую). Преломление цветных лучей при вхождении в призму оказывается разным (на рис.\,\ref{fig:p195} ход лучей показан качественно): \emph{при переходе от красного цвета к фиолетовому (при увеличении частоты света) отклонение цветного луча увеличивается}.

Зависимость показателя преломления среды от частоты света имеет вид\footnote{Эта зависимость (<<чем больше~$\nu$, тем больше~$n$>>) работает на очень малых изменениях частоты, но ей можно пользоваться для качественных ответов: например, показатель преломления для синего цвета больше показателя преломления для зеленого и~т.\,д.}:
\begin{equation}
    \label{24disp_e1}
    n\sim \nu.
\end{equation}







 
\end{document}